% ID: SEL-D0074-OA-E03
% ESTADO: cotejada. Segunda pasada visual y borrador compilado revisados.
% Fuente: Selectividad/sel_2008_matematicas_aplicadas/Mat CCSS 2 - Examen y criterios de correción.pdf
% SHA-256: 6439a5956271f683dc45ba870695a7399db4412c9b0eaf1a35f908a3857ae648
% Página: 1; opción A; ejercicio original 3.

\noindent\underline{Parte I}

El examen de Matemáticas de un alumno consta de dos ejercicios. La probabilidad de que resuelva el primero es del $30\%$, la de que resuelva ambos es del $10\%$, y la de que no resuelva ninguno es del $35\%$. Calcule las probabilidades de los siguientes sucesos:
\begin{enumerate}[label=\alph*)]
\item \textbf{(1 punto)} Que el alumno resuelva el segundo ejercicio.
\item \textbf{(1 punto)} Que resuelva el segundo ejercicio, sabiendo que no ha resuelto el primero.
\end{enumerate}
\par\noindent\begin{minipage}{\linewidth}
\noindent\underline{Parte II}

La longitud de los cables de los auriculares que fabrica una empresa es una variable aleatoria que sigue una ley Normal con desviación típica $4.5$ cm. Para estimar la longitud media se han medido los cables de una muestra aleatoria de 9 auriculares y se han obtenido las siguientes longitudes, en cm:
\[
205,\quad198,\quad202,\quad204,\quad197,\quad195,\quad196,\quad201,\quad202.
\]
\begin{enumerate}[label=\alph*)]
\item \textbf{(1 punto)} Halle un intervalo de confianza, al $97\%$, para la longitud media de los cables.
\item \textbf{(1 punto)} Determine el tamaño mínimo que debe tener una muestra de estos auriculares para que el error de estimación de la longitud media sea inferior a 1 cm, con el mismo nivel de confianza del apartado anterior.
\end{enumerate}
\end{minipage}\par
